%% ---------------------------------------------------------------------------
%% intro.tex
%%
%% Introduction
%%
%% $Id: intro.tex 1477 2010-07-28 21:34:43Z palvarado $
%% ---------------------------------------------------------------------------

\chapter{Introducción}
\label{chp:intro}

La Christian-Albrechts-Universit¨at zu Kiel conocida como CAU y Universidad de Kiel,
referente en investigación tecnológica en Alemania, ha iniciado una investigación para
diseñar una arquitectura de FPGA propia. A diferencia de las soluciones comerciales de
caja negra, este desarrollo busca ofrecer una plataforma competitiva para fines académicos
y experimentales.
El proyecto ha captado el interés de TSMC, el fabricante de semiconductores más grande
del mundo, que ha manifestado su intención de apoyar al centro educativo mediante el
acceso a nodos de fabricación avanzados. Esta sinergia se extiende a Centroamérica gracias
a los lazos históricos de cooperación entre la Universidad de Kiel y el Instituto Tecnológico
de Costa Rica, permitiendo una transferencia de conocimiento de alto nivel en el área de
sistemas digitales complejos.
La investigación se centra en el concepto de la caracterización de una FPGA embebida
en lugar de la fabricación de un IC desde cero. Su desarrollo consta del diseño de una
eFPGA en conjunto de un sistema para la inyección de flujo de bits para su correcta
programación. La esencia del proyecto es hardware reconfigurable dentro de otro para
su caracterización y desarrollo, la implementación de un filtro digital parametrizable se
requiere para términos del comportamiento de la eFPGA, entre las métricas a estudiar
se encuentra el consumo de potencia, las rutas críticas y la eficiencia de las tablas de
consulta.

\section{Planteamiento del problema}

\subsection{Estado Actual}
 Aunque existen arquitecturas de FPGA de desarrollo académico, los mismos no cuentan
con la capacidad de la flexibilidad de una FPGA comercial debido a características como
la cantidad de tablas de consulta, potencia y velocidad de procesamiento, por lo que su 
desarrollo se ha quedado estancado en su propia academia. Las FPGAs comerciales
requieren licencias para su uso además del elevado precio que significa la compra de
sufientes para impartir un curso. Además, existe un interés de TSMC para su dessarrollo,
sin embargo, aún no se cuenta con un modelo caracterizado ni validado.

\subsection{Definición del Problema}

La dependencia de arquitecturas de hardware cerrado limita la capacidad de los centros
de investigación para innovar en el diseño de semiconductores. En la colaboración CAU
TEC, el problema radica en la falta de una metodología validada para caracterizar el
rendimiento, área y consumo de una eFPGA. Sin una implementación previa en una
FPGA física que actúe como entorno de prueba para aplicaciones como filtros digitales,
no es posible asegurar que el diseño final sea competitivo para su futura fabricación con
TSMC.

\section{Estructura del Documento}

En el \textbf{capítulo~\ref{ch:marco}} se presenta un informe de la investigación del 
estado del arte actual del problema, así como una descripción de las partes y 
características del proyecto, Se profundiza en la arquitectura básica de una FPGA 
y en el método para cargarle el bitstream.

\setlength{\parskip}{2.5mm}

En el \textbf{capítulo~\ref{ch:solucion}} se presenta las posible solución al problema 
donde analizamos las etapas a tomar en cuenta para la validación de una arquitectura de 
lógica reconfigurable, se describe además cada parte del proyecto y los resultados que 
se esperan.

\setlength{\parskip}{2.5mm}

En el \textbf{capítulo~\ref{ch:HyO}}  se expone la hipótesis planteada, además, se 
presenta el objetivo general para este trabajo, así como, los objetivos específicos 
del mismo. Se muestran metas por objetivo propuesto. Por último, 
se resaltan los alcances y limitaciones que presenta este trabajo.

\setlength{\parskip}{2.5mm}

En el \textbf{capítulo~\ref{ch:Experimento}} se presenta el posible experimento 
preliminar que detalla la estrategia del proyecto y el enfoque necesario prueba el 
sistema y el flujo de herramientas a través de un experimento bien diseñado y riguroso.

\setlength{\parskip}{2.5mm}

En el \textbf{capítulo~\ref{ch:Crono}}  se presentan las tareas necesarias a llevar a 
cabo a lo largo de este trabajo, así como el diagrama de Gantt para poder completar con las 
tareas propuestas.

\setlength{\parskip}{2.5mm}

%%% Local Variables: 
%%% mode: latex
%%% TeX-master: "main"
%%% End: 
